\subsubsection{Topological sort (Kahn + DFS)}
\begin{lstlisting}[style=mystyle, language=C++]
// TC: O(N + M)
// usa: orden lineal de un DAG
#include <bits/stdc++.h>
using namespace std;

vector<int> topoKahn(int n, const vector<vector<int>>& adj){
	vector<int> indegree(n, 0);
	for(int i=0 ; i<n ; i++) for(auto adjN : adj[i]) indegree[adjN]++;
	queue<int> q;
	for(int i=0 ; i<n ; i++) if(indegree[i] == 0) q.push(i);
	vector<int> order;
	while(!q.empty()){
		int node = q.front(); q.pop();
		order.push_back(node);
		for(auto adjN : adj[node]){
			if(--indegree[adjN] == 0) q.push(adjN);
		}
	}
	return order; // size < n if cycle
}

void dfsVisit(int node, const vector<vector<int>>& adj, vector<int>& state, vector<int>& order){
	state[node] = 1;
	for(auto adjN : adj[node]){
		// if(state[adjN] == 1) -> ciclo: back edge a ancestro en la pila de recursion
		if(state[adjN] == 0)
			dfsVisit(adjN, adj, state, order);
	}
	state[node] = 2;
	order.push_back(node);
}

vector<int> topoDFS(int n, const vector<vector<int>>& adj){
	vector<int> order;
	vector<int> state(n, 0); // 0=white, 1=gray, 2=black
	for(int i=0 ; i<n ; i++) 
		if(state[i] == 0) 
			dfsVisit(i, adj, state, order);
	reverse(order.begin(), order.end());
	return order;
}

int main(){
	vector<vector<int>> adj = {{1,2}, {3}, {3}, {}};
	auto order = topoKahn(4, adj);
	for(int node : order) cout << node << " ";  // 0 1 2 3
	cout << "\n";
	return 0;
}
\end{lstlisting}

\subsubsection{DP en DAG (camino mas largo / corto con relajacion)}
\begin{lstlisting}[style=mystyle, language=C++]
// TC: O(N + M)
// usa: Procesar en orden topologico relaja cada arista EXACTAMENTE 1 vez (Bellman-Ford lineal).
// Sirve para camino mas largo/corto con pesos arbitrarios (positivos o negativos) y conteo de caminos.
const long long INF = 1e18;

// Retorna {dist, parent}. dist[i] = -INF si i es inalcanzable desde s.
pair<vector<long long>, vector<int>> dagLongestPath(int n, const vector<vector<pair<int,long long>>>& adj, const vector<int>& topo, int s){
	vector<long long> dist(n, -INF);
	vector<int> parent(n, -1);
	dist[s] = 0; // o 1 si se cuentan ciudades/vertices

	for(auto node : topo){
		if(dist[node] == -INF) continue; // Inalcanzable desde s
		for(auto [adjN, adjW] : adj[node]){
			if(dist[node] + adjW > dist[adjN]){
				dist[adjN] = dist[node] + adjW;
				parent[adjN] = node; // Para reconstruir la ruta
			}
		}
	}
	return {dist, parent};
}

// Reconstruir camino de s a t (retorna vacio si t es inalcanzable)
vector<int> getPath(int s, int t, const vector<int>& parent){
	if(s != t && parent[t] == -1) return {};
	vector<int> path;
	for(int cur = t; cur != -1; cur = parent[cur]) path.push_back(cur);
	reverse(path.begin(), path.end());
	if(path.empty() || path[0] != s) return {};
	return path;
}

// Conteo de caminos s -> t (modulo MOD):
// vector<long long> ways(n, 0); ways[s] = 1;
// for(auto node : topo) if(ways[node]) for(auto [adjN, adjW] : adj[node]) ways[adjN] = (ways[adjN] + ways[node]) % MOD;
\end{lstlisting}
